\documentclass{article}
\usepackage{amsmath,amssymb,amsthm}
\usepackage[margin=1in]{geometry}
\usepackage{url}
\newtheorem{lemma}{Lemma}
\newtheorem{proposition}{Proposition}
\newcommand{\Norm}{\mathrm N}
\newcommand{\C}{\mathbb C}
\begin{document}

\title{Complete means of the exceptional additive Gram blocks}
\author{}
\date{October 8, 2026}
\maketitle

This note records arithmetic deductions from the pinned upstream revision
\texttt{adc7f1241}\allowbreak\texttt{b42e322a645}\allowbreak
\texttt{1854ab7e4b4c}\allowbreak\texttt{146bf78a}.
It does not assert a new zero-free boundary.
The local finite-field formulas below have direct proofs.
Run their exact checks with
\begin{quote}
\texttt{uv run research/gram\_bulk.py}.
\end{quote}
The Lean file \path{lean/GramBulkClassification.lean} records the vanishing
of the pulled-back coefficient when the dilation meets a common prime.

\section*{The coefficient being averaged}

Let $H$ be the fixed modulus called \texttt{jointFixedModulus} upstream.
For the common ideal $C$, frequency $k\ne0$, and dilation $d$, the mean is
\[
 \mu_{C,d,k}=\frac1{q_{\mathfrak r}^2}
 \sum_{x,y\bmod\mathfrak r}
 \operatorname{canonicalJoint}(C,k,dx,dy),
\]
where $\mathfrak r$ is any period of the pulled-back coefficient.
The multiplication by $d$ occurs before averaging.
If $d$ is a unit modulo $\mathfrak r$, multiplication permutes its residue
classes and $\mu_{C,d,k}=\mu_{C,1,k}$.
If a prime of $C$ divides $d$, the mean is zero.
Indeed \texttt{localExtension} is defined to be zero when both columns vanish
at that prime.
If a prime outside the fixed modulus divides both $d$ and $k$, the moving
numerator character supplies a zero mask and the coefficient is again zero.
The calibration masks also force zero if $C$ or $d$ meets its excluded set.
These facts must be checked before replacing a pulled-back mean by an
undilated mean.

The exact Poisson bulk is
\[
 \operatorname{Bulk}_{C,d,k}
 =\frac43\mu_{C,d,k}N^2\int_{\C^2}W_{C,d,k}.
\]
It is the zero dual frequency.
It is not the lattice sum with the arithmetic coefficient replaced by its
complete mean.
Subtracting that latter lattice sum leaves an additional constant-coefficient
lattice discrepancy.

\section*{Good primes outside the fixed modulus}

Fix a finite field $F$ of cardinality $q$ and a multiplicative character
$\chi$ of exact order six, extended by zero at zero.
Write $c=v_p(C)$ and $e=v_p(k)$.
For $e=0$, define $\psi_e(u)=1$ for every $u\in F$.
For $e>0$, define $\psi_e(u)=\chi(u)^e$, with value zero at $u=0$.
For $c=0$, put $L_c(u,v;k)=1$.
For $c>0$, put
\[
 L_c(u,v;k)=
 \begin{cases}
 0,&u=v=0,\\
 q^{c-1}\displaystyle\sum_{vx-uy=k}
        \chi(x)^c\overline{\chi(y)^c},&\text{otherwise}.
 \end{cases}
\]
When $e=0$, $k$ is any nonzero element of $F$.
When $e>0$, its residue is zero.
This is the local factor in \texttt{GramLocalExtension.lean}, multiplied by
the two moving numerator factors in \texttt{GramJoint.lean}.
Define
\[
 m_q(c,e;d)=q^{-2}\sum_{u,v\in F}
       L_c(du,dv;k)\psi_e(du)\overline{\psi_e(dv)}.
\]

\begin{proposition}
If $d\ne0$ in $F$, then
\[
 m_q(c,e;d)=
 \begin{cases}
 1,&c=e=0,\\
 0,&6\nmid c+e,\\
 (1-q^{-1})^2,&c=0, e>0, 6\mid e,\\
 q^c(1-q^{-1})^2,&c>0, e=0, 6\mid c,\\
 q^c(1-q^{-1})^3,&c>0, e>0, 6\mid c+e.
 \end{cases}
\]
If $d=0$, the mean is zero unless $c=e=0$, in which case it is one.
\end{proposition}

\begin{proof}
A nonzero dilation permutes $F$ in both columns, so it suffices to take $d=1$.
If $c=0$ and $e>0$, character orthogonality gives
\[
 q^{-2}\left|\sum_{u\in F}\chi(u)^e\right|^2
 =\begin{cases}(1-q^{-1})^2,&6\mid e,\\0,&6\nmid e.\end{cases}
\]
Suppose $c>0$ and $e>0$.
The numerator factors kill both one-zero branches.
For unit columns and $k=0$, the local correlation is
$(q-1)\chi(u/v)^c$.
The weighted coefficient is therefore
$q^{c-1}(q-1)\chi(u/v)^{c+e}$.
The sum over the two unit columns is $(q-1)^2$ if $6\mid c+e$ and zero
otherwise.
After division by $q^2$, this proves the last branch.

Suppose $c>0$ and $e=0$, so $k\ne0$.
For unit columns, the correlation is
\[
 \chi(u/v)^c
 \begin{cases}q-2,&6\mid c,\\-1,&6\nmid c.\end{cases}
\]
The two one-zero branches are each $q-1$ if $6\mid c$ and zero otherwise.
When $6\mid c$, their total before multiplication by $q^{c-1}$ is
$(q-1)^2(q-2)+2(q-1)^2=q(q-1)^2$.
When $6\nmid c$, the unit-column sum vanishes by orthogonality and both
one-zero branches vanish.
This proves the fourth branch.
Finally, if $d=0$ and $c>0$, the definition of $L_c$ kills the coefficient.
If $d=0$, $c=0$, and $e>0$, the numerator zero mask kills it.
\end{proof}

Thus a nonzero complete mean requires the combined ideal $C(k)$ to have
valuations divisible by six away from the fixed support.
The existing exceptional count imposes this condition only on $(k)$ at primes
outside $C$ and the fixed support.
The complete mean imposes the stronger combined condition.
All common-part exponents are included in this proposition.
No squarefree-column hypothesis occurs.

\section*{The common-part sum}

\begin{lemma}
For every integer $h>0$ divisible by six,
\[
 \sum_{c=0}^h q^{-2c}m_q(c,h-c;1)
 =(1-q^{-1})(1-q^{-2}).
\]
\end{lemma}

\begin{proof}
Put $r=q^{-1}$ and factor out $(1-r)^2$.
The endpoint contributions are $1$ and $r^h$.
The interior contributions sum to
$(1-r)\sum_{c=1}^{h-1}r^c=r-r^h$.
The resulting factor is $1+r$.
\end{proof}

This identity is positive and independent of $h$.
Consequently these good-prime common-part branches do not cancel each other.
If the complete signed M\"obius sum can be extended to all dilation ideals,
the primes outside $C(k)$ and the fixed support contribute
$\prod_p(1-q_p^{-2})$.
Combining this product with the preceding identity formally gives
\[
 \zeta_{\mathrm{good}}(2)^{-1}
 \prod_{p\mid C(k),\ p\ \mathrm{good}}(1-q_p^{-1}).
\]
The extension of the finite divisor pool to this infinite Euler product needs
a tail estimate.
This note does not identify the finite divisor-pool sum with that product.

\section*{Good excluded primes}

Let $p$ lie in the excluded set $S$, away from the primes over $2,3$.
The calibration zero masks require $c=0$ and $d$ a unit at $p$.
The generator $b_*$ has valuation one at $p$.
Hence the local character in the first column is
\[
 \chi_p^{1+e}\overline{\xi_p},
\]
and the second column has its conjugate.
The local complete mean is zero unless $\chi_p^{1+e}=\xi_p$.
When this identity holds, its value is $(1-q_p^{-1})^2$.
For the quadratic calibration $\xi_p=\chi_p^3$, the necessary and sufficient
condition at this prime is $e\equiv2\pmod6$.
This local assertion follows from character orthogonality and the upstream
definitions of \texttt{gramPeriodicMonoid} and \texttt{jointExtension}.
Its assembly with the remaining fixed phases is not yet formalized.

\section*{The remaining fixed phase}

The primary restriction, physical ray restriction, supplementary phases at
$2,3$, and the reciprocity phase remain in a finite complete mean.
They cannot be discarded.
Their exact coefficient is the function \texttt{jointFixed} upstream,
with the $S$-supported moving numerator factors retained.
A global formula must show how that finite mean depends on the residue class
of $C$, the unit and bad-prime components of $k$, and the ray sector.
The local product above is a proved factor only at primes coprime to that
fixed modulus.

\begin{lemma}
If $m$ is supported and primary, and $4\mid n-m$, then
\[
 \operatorname{jointFixed}(C,k;n,m)=f(n)\overline{f(m)},
\]
where
\[
 \begin{split}
 f(n)={}&\operatorname{primaryCoefficient}(Cn)\\
 &\cdot\operatorname{numeratorBadTwist}(u,a,b,r;n).
 \end{split}
\]
\end{lemma}

\begin{proof}
The diagonal reciprocity identity is $\mathcal R(m,m)=\chi_m(-1)$.
For primary prime generators this is the upstream identity in
\texttt{Detector/PrimePhase.lean}.
It extends to every supported primary $m$ by prime factorization.
Indeed the reciprocity phase is multiplicative in both arguments, symmetric,
and takes values $\pm1$ on supported arguments, so
$\mathcal R(ab,ab)=\mathcal R(a,a)\mathcal R(b,b)$.
The denominator character is also multiplicative in the denominator ideal.

Since $n$ and $m$ have the same ray modulo $4$,
$\mathcal R(n,m)=\mathcal R(m,m)=\chi_m(-1)$.
Also
\[
 \operatorname{numeratorBadTwist}(-u;m)
 =\chi_m(-1)\operatorname{numeratorBadTwist}(u;m).
\]
The extra two factors in \texttt{jointFixed} multiply to
$\chi_m(-1)\overline{\chi_m(-1)}=1$.
\end{proof}

The general diagonal identity and this rank-one lemma are compiled in
\texttt{GramBulkClassification.lean} with only the standard axioms
\texttt{propext}, \texttt{Classical.choice}, and \texttt{Quot.sound}.
To apply the lemma without explicit column hypotheses, one must use the
original masks.
The numerator unit supplement kills unsupported columns.
For primary $C$, the primary restriction on $Cm$ implies the one on $m$.
Since supported $C$ is invertible modulo $4$, the two ray restrictions on
$Cn,Cm$ imply $4\mid n-m$.
That bridge is not yet formalized here.

\section*{An explicit nonzero arithmetic mean}

Assume $S$ contains the unique primes over $2$ and $3$, and let
$S_{\mathrm g}$ be its other primes.
Write the squarefree calibration generator as
\[
 b_*=u_b\lambda\,2\prod_{p\in S_{\mathrm g}}\pi_p,
\]
with primary good generators $\pi_p$.
Every sixth-torsion residue character at a good prime is a power of its
canonical exact-order-six character.
Thus choose $1\le a_p\le5$ such that $\xi_p=\chi_p^{a_p}$, and put
\[
 k_0=u_b^{-1}\lambda^5\,2^5
       \prod_{p\in S_{\mathrm g}}\pi_p^{a_p-1}.
\]
Then $b_*k_0=\lambda^6\,2^6\prod_p\pi_p^{a_p}$.

Fix a unit ray $\sigma$ modulo the ideal $(4)$.
For $C=d=1$, the single-column coefficient after the rank-one reduction is
\[
 f(n)=\tau^{-1}\xi(n)^{-1}\chi_n(b_*k_0)
\]
on primary columns in ray $\sigma$ coprime to $S$, and zero otherwise.
On those columns it is a constant $\gamma_\sigma$ of norm one.
To see this, the sixth powers at the bad primes contribute one, with their
unit masks retained.
At a good excluded prime, reciprocity gives
$\chi_n(\pi_p)^{a_p}
=\mathcal R(\pi_p,n)^{a_p}\chi_p(n)^{a_p}$.
The character cancels $\xi_p(n)$, and the phase is constant in ray $\sigma$.
The calibration character at $\lambda$ is fixed by
$n\equiv1\pmod{\lambda^2}$.
The calibration character at $2$ is fixed by $n\bmod4$.

The ideals $(\lambda^2)$ and $(4)$ are coprime and have norms $9$ and $16$.
CRT therefore gives the active residue density
\[
 \delta_\sigma=\frac1{144}
     \prod_{p\in S_{\mathrm g}}(1-q_p^{-1}).
\]
Consequently
\[
 \mu_{1,1,k_0}=\delta_\sigma^2>0.
\]
For a supported element $a$ coprime to $S$, multiplication of the frequency
by $a^6$ retains the unit masks at its primes and gives
\[
 \mu_{1,1,k_0a^6}
 =\delta_\sigma^2\prod_{p\mid(a)}(1-q_p^{-1})^2.
\]
The finite-field character spectrum and CRT density in this argument have
not yet been formalized in Lean.
The formula is a mathematical deduction using those standard finite-group
facts and the rank-one identity.
It proves nonvanishing of the arithmetic mean.
It does not prove nonvanishing or positivity of the weighted bulk, because
the smooth integral and subsequent frequency sum can cancel.

\section*{The two finite pools}

Write $Y$ for the physical column scale, including any preceding rescaling.
The upstream window is the unmasked set
\[
 F_Y=\{I\text{ supported}:W(q_I/Y)\ne0\}.
\]
It does not impose a ray or calibration restriction.
The actual common pool is $\{\gcd(I,J):I,J\in F_Y\}$.
The dilation pool at $C$ is the union of all ideal divisors of the residual
window $F_{Y/q_C}$.
These are the definitions in \texttt{GramFiniteGcd.lean:16},
\texttt{LowColumnBounds.lean:13}, and \texttt{MobiusRegroup.lean:11}.

Choose $0<a'<b'$ such that $W$ never vanishes on $[a',b']$, and put
\[
 L_0=\max\left(1,\left(\frac{12}{\sqrt{b'}-\sqrt{a'}}\right)^2\right).
\]
If $q_C\le Y/L_0$, the interval
$[\sqrt{a'Y/q_C},\sqrt{b'Y/q_C}]$ has length at least $12$.
It contains $j,j+6$ with $j\equiv1\pmod6$ and $j>0$.
Both integer ideals are supported and they are coprime.
Thus $C(j),C(j+6)\in F_Y$ and their gcd is exactly $C$.
Every supported $C$ up to $Y/L_0$ belongs to the actual common pool.

Similarly, for $N=Y/q_C$, every supported $D$ with $q_D\le N/L_0$
divides a column $D(j)$ in $F_N$.
Here one only needs $j\equiv1\pmod6$ in the interval
$[\sqrt{a'N/q_D},\sqrt{b'N/q_D}]$.
Extra good excluded primes dividing $j$ cause no problem because the pools
are unmasked.

The elementary ideal count $\#\{D:q_D\le T\}\ll T$ gives
\[
 \sum_{q_D>L}q_D^{-2}\ll L^{-1}\qquad(L\ge1).
\]
Consequently, for any specified set of allowed dilation primes, the actual
finite signed sum $\sum_D\mu(D)q_D^{-2}$ differs from its infinite Euler
product by $O_W(q_C/Y)$.
For $Y/q_C<L_0$, the same bound follows by bounding both sums absolutely.
This argument concerns the signed factor only after its nonunit dilation
terms have been proved to vanish.
It does not replace a raw Poisson remainder by a centered lattice sum.

\section*{Assembly at one combined sixth-power frequency}

Keep $k_0$ and the unit ray from the explicit nonzero family above.
Let $A$ be an ideal coprime to $S$, let $a$ be its multiplicative primary
generator, and put $H=k_0a^6$.
Only common ideals $C\mid A^6$ can have nonzero mean at this combined
frequency.
Set $k=H/\operatorname{primaryGenerator}(C)$.
At an excluded good prime $k$ retains exponent $a_p-1$.
At a prime outside $S$, write $h=6v_p(A)$ and $c=v_p(C)$.
The complete mean is
\[
 \mu_{C,1,k}=\delta_\sigma^2
   \prod_{p\mid A}m_{q_p}(c,h-c;1).
\]
The fixed density is independent of $C$.
The original ray condition becomes the unit ray $\sigma/C$ on the residual
column, which has the same density.
The corresponding constant phase cancels with its conjugate.
If $D$ meets $S$ or $A$, its pulled-back coefficient vanishes.
Otherwise it is a unit at every relevant period prime and its mean is the
displayed one.

Put $P=Y^2/Q$.
The physical frequency scale is exactly
\[
 \operatorname{frequencyScale}(C,k,Y,Q)
   =q_C\Norm(k)/P=\Norm(H)/P.
\]
The profile integral after lattice scaling therefore depends on $H$ and $P$,
and not on $C$ or $D$.
Write it as $J_v(\Norm(H)/P)$.
The zero dual frequency has prefactor
\[
 \frac{Q}{Y^3}\frac43\left(\frac{Y}{q_Cq_D}\right)^2
 =\frac43\frac QY q_C^{-2}q_D^{-2}.
\]
This is the exact Poisson bulk, with the standard Lebesgue integral on
$\mathbb C^2$.

Let $\zeta_S(s)=\prod_{p\notin S}(1-q_p^{-s})^{-1}$.
The infinite signed dilation factor is
\[
 E_A=\zeta_S(2)^{-1}\prod_{p\mid A}(1-q_p^{-2})^{-1}.
\]
The local common-part identity then gives
\[
 E_A\sum_{C\mid A^6}q_C^{-2}\mu_{C,1,k}
 =\delta_\sigma^2\zeta_S(2)^{-1} f(A),\qquad
 f(A)=\prod_{p\mid A}(1-q_p^{-1}).
\]

Both finite-pool errors can be bounded before summing over $A$.
Since $0\le\mu_{C,1,k}\le\delta_\sigma^2q_C$,
\[
 \sum_{C\mid A^6}q_C^{-1}\mu_{C,1,k}
 \le\delta_\sigma^2\tau(A^6).
\]
The finite dilation error at $C$ is $O_W(q_C/Y)$.
Any omitted common ideal has $q_C>Y/L_0$, and $0<E_A\le1$.
It follows that the actual arithmetic factor $T_Y(A)$ satisfies
\[
 \left|T_Y(A)-\delta_\sigma^2\zeta_S(2)^{-1}f(A)\right|
 \ll_{S,W}\frac{\tau(A^6)}Y.
\]
The common profile $J_v(\Norm(H)/P)$ can be factored out before applying
this bound.
The pool saturation and this error estimate are paper-level proofs from the
actual finite definitions, not yet Lean theorems.

\section*{The weighted family asymptotic}

Assume $W\in C_c^\infty((0,\infty))$ is nonzero and $U$ is a Schwartz
function on $\mathbb R$.
Here $U(r)$ is used as a radial function $U(|z|^2)$ on $\mathbb C$.
For $s\in\mathbb C$, write
$M_W(s)=\int_0^\infty W(r)r^s\,dr/r$.
With $F_U=\operatorname{paperRadialFourier}(U)$, the common integral is
\[
 J_v(T)=\pi^2\int_0^\infty\int_0^\infty
 W(r_1)\overline{W(r_2)}
 r_1^{-1+iv}r_2^{-1-iv}F_U(T/(r_1r_2))\,dr_1\,dr_2.
\]
The factor $\pi^2$ is from polar coordinates in the two complex variables.
For real $v$ and $\alpha>0$,
\[
 M_{J_v}(\alpha)
 =\pi^2|M_W(\alpha+iv)|^2 M_{F_U}(\alpha).
\]
The profile and its $T$ derivative decay rapidly at infinity, uniformly in
real $v$, because the oscillatory powers have modulus one.

The Eisenstein ideal count coprime to $S$ is
\[
 A_S(T)=c_ST+O_S(\sqrt T+1),\qquad
 c_S=\frac{\pi}{3\sqrt3}\prod_{p\in S}(1-q_p^{-1}).
\]
Indeed there are six generators for each nonzero ideal and the Eisenstein
lattice has covolume $\sqrt3/2$.
Finite inclusion-exclusion gives the displayed coprimality factors.
Since $f(A)=\sum_{D\mid A}\mu(D)/q_D$, convolution gives
\[
 \sum_{q_A\le T,\ (A,S)=1}f(A)
 =\frac{c_S}{\zeta_S(2)}T+O_S(\sqrt T+1).
\]
The error uses the convergent series $\sum_Dq_D^{-3/2}$ and the tail
$\sum_{q_D>T}q_D^{-2}\ll T^{-1}$.

Put $X=(P/\Norm(k_0))^{1/6}$.
Partial summation gives
\[
 \sum_{(A,S)=1} f(A)J_v((q_A/X)^6)
 =\frac{c_S}{6\zeta_S(2)}X M_{J_v}(1/6)
   +O_{S,W,U}(\sqrt X+1).
\]
This error is uniform for all real $v$.
For example, it is bounded by a constant times
\[
 \sqrt X\int_0^\infty\sqrt u\left|\frac d{du}J_v(u^6)\right|du
 +\int_0^\infty\left|\frac d{du}J_v(u^6)\right|du.
\]
Both integrals are uniformly finite.
The divisor estimate $\tau(A^6)\ll_\eta q_A^\eta$ and rapid profile decay
bound the aggregate finite-pool error by
$O_{S,W,U,\eta}((X^{1+\eta}+1)/Y)$.

Thus the exact zero-dual-frequency contribution from this one specified
combined-frequency family is, for $P,Y\ge1$ and any $\eta>0$,
\[
 \begin{split}
 \mathcal B_{k_0}(P,Y,v)
 ={}&\frac{2\pi^2}{9}
 \frac{\delta_\sigma^2c_S}{\zeta_S(2)^2}
 \Norm(k_0)^{-1/6}\frac QY P^{1/6}
 |M_W(1/6+iv)|^2 M_{F_U}(1/6)\\
 &+O_{S,W,U,\eta}\left(
 \frac QY\left[P^{1/12}+1+
 \frac{P^{(1+\eta)/6}+1}{Y}\right]\right).
 \end{split}
\]
Each ideal $A$ supplies the single element $H=k_0a^6$.
There is no additional factor of six in this chosen family.
Other unit or bad-support frequency families have not been included.
This theorem is a mathematical derivation using the mean classification,
pool arguments, and lattice count above.
Its CRT and spectral ingredients remain unformalized in Lean.

\section*{Fourier normalization and the radial obstruction}

Upstream \texttt{SexticRadialPoisson.lean:304} defines
\[
 \operatorname{paperFourier}(g)(u)
 =\frac2{\sqrt3}\widehat g\left(\frac{2i}{\sqrt3}\overline u\right),
\]
where the standard Fourier transform has phase $e^{-2\pi i x\cdot\xi}$.
For a radial function this gives
\[
 F_U(t)=\frac{2\pi}{\sqrt3}\int_0^\infty
 U(r)J_0\left(\frac{4\pi}{\sqrt3}\sqrt{rt}\right)dr.
\]
At $\alpha=1/6$, the absolute convergence of the Bessel Mellin integral
justifies Fubini and gives\footnote{The identity
$\int_0^\infty x^{2\alpha-1}J_0(x)dx
=2^{2\alpha-1}\Gamma(\alpha)/\Gamma(1-\alpha)$ is
DLMF 10.22.43, \url{https://dlmf.nist.gov/10.22.E43}.}
\[
 M_{F_U}(\alpha)
 =\frac{2\pi}{\sqrt3}
 \left(\frac3{4\pi^2}\right)^\alpha
 \frac{\Gamma(\alpha)}{\Gamma(1-\alpha)}M_U(1-\alpha).
\]
Near infinity its absolute integrand is
$O(x^{2\alpha-3/2})$, which is integrable at $\alpha=1/6$.
Near zero it is $O(x^{2\alpha-1})$.
Hence if $U$ is real, nonnegative on $(0,\infty)$, and nonzero there,
$M_{F_U}(1/6)>0$.
For such $U$, if also $M_W(1/6+iv)\ne0$, the displayed leading coefficient
of this family is positive.
For fixed $v$ with this nonvanishing property, it is a genuine leading
asymptotic as $P\to\infty$ and $P^{\eta/6}/Y\to0$.
The absolute error is uniform in $v$.
A relative asymptotic in a growing frequency range requires an explicit
lower bound on $|M_W(1/6+iv)|$ compared with that error.

The smooth functional has the single-window Mellin factor from the earlier
radial projection analysis:
\[
 |M_W(1/6+iv)|^2M_{F_U}(1/6).
\]
For $D=r\,d/dr$, replacing $W$ by $(D+1/6+iv)W$ annihilates its first
Mellin factor by integration by parts.
For the physical two-window kernel
$W_0DW_1-DW_0W_1+(1/6)W_0W_1$, the source multiplier is $7/6-w-z$.
It vanishes at $(w,z)=(1,1/6)$.
In that source convention the principal $W_1$ Mellin value is at $1$,
and the principal radial Fourier value of $W_0$ is at $1/6$.
The $U$ here is the Gram energy test and is not assumed to be that physical
$W_0$.
Finite coupled window combinations require their cross Gram terms before
this single-window calculation can be identified with the full coupled
Mellin functional.
It does not prove that this family is the entire exceptional bulk.
Other frequency families may cancel a signed subfamily contribution.
Nor does it compare the bulk with the centered remainder without the exact
Poisson split.
There is no new zero-free boundary conclusion here.

\section*{Checks and remaining obligations}

The script computes the defining finite-field correlation directly using
integer pairs in $\mathbb Z[\omega]$.
It uses no floating-point roots of unity.
At each of $q=7,13,19$, it checks all $0\le c,e\le6$ and dilations $0,1,2$,
giving $441$ exact mean checks.
It also checks the common-part sum at $h=6,12,18$.
For example, at $q=7$,
\[
 m_7(1,5;1)=216/49,\qquad m_7(1,0;1)=0,\qquad
 m_7(6,0;1)=86436.
\]

The remaining formalization obligations are CRT assembly, the active-column
bridge for the rank-one identity, the finite-field character spectrum,
the two pool saturation lemmas, and the weighted ideal count.
The remaining mathematical classification problem is the other unit and
bad-support families and their signed aggregate.
Only after that classification can the full exceptional bulk be compared
with the detector's principal contribution.

\end{document}
